\subsection{替换}

的多项式。设 E 是 A 上的含单位元结合代数，且 \(\mathbf{x} = (x_{i})_{i \in I}\) 是 E 中两两可交换的元素族。设 \(\mathbf{X} = (\mathbf{X}_{i})_{i \in I}\) 是一族未定元。由 III，第 449 页命题 7，存在唯一的从 A{[}X{]} 到 E 中的保单位同态 f，使得 \(f(\mathbf{X}_{i}) = x_{i}\) （对每个 \(i \in I\) ）。A{[}X{]} 的元素 u 在 f 下的像记为 \(u(\mathbf{x})\) ，称为在 u 中以 \(x_{i}\) 代替 \(X_{i}\) 所得的 E 的元素，也称为 u 在 \(X_{i} = x_{i}\) 处的值。特别地， \(u = u((\mathbf{X}_{i})_{i \in I})\) 。若 \(I = \{1, \ldots, n\}\) ，我们记 \(u(x_{1}, \ldots, x_{n})\) 以代替 \(u((x_{i})_{i \in I})\) 。更一般地，如果 M 是一个 A-模，并且 v 是

\[
\mathbf {M} [ (\mathbf {X} _ {i}) _ {i \in I} ] = \mathbf {M} \otimes_ {\mathbf {A}} \mathbf {A} [ (\mathbf {X} _ {i}) _ {i \in I} ],
\]

的一个元素，我们把 \(v\) 在 \(\mathbf{M} \otimes_{\mathbf{A}} \mathbf{E} = \mathbf{M}_{(\mathbf{E})}\) 中在映射 \(1_{\mathbf{M}} \otimes f\) 下的像记为 \(v(\mathbf{x})\) 。

如果同态 \(u \mapsto u(\mathbf{x})\) （从 \(\mathbf{A}[\mathbf{X}]\) 到 \(\mathbf{E}\) 中的同态）是单射，我们就说族 \(\mathbf{x}\) 在 \(\mathbf{A}\) 上是代数自由的，或者这些 \(x_i\) 在 \(\mathbf{A}\) 上代数无关。这也就是说，单项式 \(\mathbf{x}^\nu\) ( \(\nu \in \mathbf{N}^{(I)}\) ) 在 \(\mathbf{A}\) 上线性无关。

如果 \(\lambda\) 是 \(E\) 到某个含单位元的结合 A-代数 \(E'\) 中的保单位同态，则我们有

\[
\lambda \left(u \left(\left(x _ {i}\right) _ {i \in I}\right)\right) = u \left(\left(\lambda \left(x _ {i}\right) _ {i \in I}\right)\right),\tag{1}
\]

这是因为 \(\lambda \circ f\) 是 \(\mathbf{A}[\mathbf{X}]\) 到 \(\mathbf{E}'\) 中的同态，它把 \(\mathbf{X}_i\) 映为 \(\lambda(x_i)\) 。

设 \(u \in A[X]\) 。若 E 是交换的，则映射 \(x \mapsto u(x)\) （从 \(E^{I}\) 到 E 中）称为由 \(u\) （及代数 E）定义的多项式函数；我们有时把它记作 \(\tilde{u}\) （甚至就记作 \(u\) ）。

设 \(\mathbf{Y} = (\mathbf{Y}_j)_{j\in J}\) 是另一族未定元，并取 E 为多项式代数 \(\mathbf{A}[\mathbf{Y}]\) 。给定 \(u\in \mathbf{A}[\mathbf{X}]\) ，取 \(g_{i}\in \mathbf{A}[\mathbf{Y}]\) （对每个 \(i\in I\) ），并令 \(\mathbf{g} = (g_i)_{i\in I}\) ；设 \(u(\mathbf{g})\in \mathbf{A}[\mathbf{Y}]\) 为以多项式 \(g_{i}\) 代替 \(\mathbf{X}_i\) 代入 \(u\) 后所得的多项式。设 \(\mathbf{y} = (y_j)_{j\in J}\) 为某个含单位元的结合 A-代数 \(\mathbf{E}'\) 中两两可交换的一族元素；应用（1）并取 \(\lambda\) 为同态 \(g\mapsto g(\mathbf{y})\) （从 E 到 \(\mathbf{E}'\) 中），我们得到

\[
(u (\mathbf {g})) (\mathbf {y}) = u \left(\left(g _ {i} (\mathbf {y})\right)\right).\tag{2}
\]

如果 \(\mathbf{f} = (f_i)_{i\in I}\in (\mathbf{A}[(\mathbf{X}_j)_{j\in J}])^I\) 且 \(\mathbf{g} = (g_j)_{j\in J}\in (\mathbf{A}[(\mathbf{Y}_k)_{k\in K}])^J\) ，我们用 \(\mathbf{f}\circ \mathbf{g}\) 或 \(f(\mathbf{g})\) 表示多项式族 \((f_i(\mathbf{g}))_{i\in I}\in (\mathbf{A}[(\mathbf{Y}_k)_{k\in K}])^I\) 。如果我们用 \(\tilde{\mathbf{f}}\) 表示映射 \(\mathbf{x} \mapsto (f_i(\mathbf{x}))_{i \in I}\) （从 \(\mathrm{E}^{\prime J}\) 到 \(\mathrm{E}^{\prime I}\) 中）（其中 \(\mathrm{E}'\) 是一个含单位元、结合且交换的 A-代数），则关系式 (2) 蕴含

\[
(\mathbf {f} \circ \mathbf {g}) ^ {\sim} = \tilde {\mathbf {f}} \circ \tilde {\mathbf {g}}.\tag{3}
\]

如果 \(\mathbf{h} = (h_k)_{k\in K}\in (\mathbf{A}[(\mathbf{Z}_l)_{l\in L}])^{\mathbf{K}}\) ，由(2)可得：

\[
\mathbf {f} \circ (\mathbf {g} \circ \mathbf {h}) = (\mathbf {f} \circ \mathbf {g}) \circ \mathbf {h}.\tag{4}
\]

命题2。——设 \(\mathbf{a} = (a_i)_{i \in I}\) 是 A 中元素的一个族，并令 \(u \in A[X]\) 。若 v 是以 \(X_i + a_i\) 代替 \(X_i\) （对每个 \(i \in I\) ）代入后所得的多项式，则 v 的常数项等于 \(u(\mathbf{a})\) 。

v 的常数项是通过把 v 中的 \(X_{i}\) （对每个 \(i \in I\) ）替换为 0 而得到的。因此结果由 (2) 得出。

推论 1. --- 设 m 为满足 \(u \in A[X]\) 且 \(u(\mathbf{a}) = 0\) 的多项式 u 所成的理想。则 m 由多项式 \(X_{i} - a_{i}\) （对于 \(i \in I\) ）生成。

显然 \(X_{i} - a_{i} \in m\) （对每个 \(i \in I\) ）。设 \(u \in m\) ，并令 v 如命题 2 所述。由于 v 没有常数项，因此存在一个多项式族 \((\mathbf{P}_{i})_{i \in I}\) （由 A{[}X{]} 中的多项式组成，具有有限支撑），使得

\[
v (\mathbf {X}) = \sum_ {i \in I} \mathrm{X} _ {i} \cdot \mathrm{P} _ {i} (\mathbf {X}).
\]

如果在上述方程中把 \(X_{i}\) 替换为 \(X_{i}-a_{i}\) （对每个 \(i\in I\) ），我们就得到一个形如 \(u(\mathbf{X})=\sum_{i\in I}\left(\mathbf{X}_{i}-a_{i}\right)\cdot\mathbf{P}_{i}^{\prime}(\mathbf{X})\) 的关系式，由此得出推论。

推论 2. —— 设 \(\mathbf{X} = (\mathbf{X}_i)_{i \in I}\) 与 \(\mathbf{Y} = (\mathbf{Y}_i)_{i \in I}\) 是两组未定元。多项式 \(u \in \mathbf{A}[\mathbf{X}, \mathbf{Y}]\) 中使得 \(u(\mathbf{X}, \mathbf{X}) = 0\) 者所成的集合，是 \(\mathbf{A}[\mathbf{X}, \mathbf{Y}]\) 的由多项式 \(\mathbf{X}_i - \mathbf{Y}_i\) （对于 \(i \in I\) ）所生成的理想。

此推论直接由推论 1 得出：只需把 A 替换为 A{[}Y{]} ，把 \(a_{i}\) 替换为 \(Y_{i}\) ，并把 A{[}X, Y{]} 解释为以 \(X_{i}\) 为未定元、系数在 A{[}Y{]} 中的多项式环。

命题3. --- 设 \(u \in \mathbf{A}[\mathbf{X}]\) ，并令 \(\mathbf{X} \cdot \mathbf{Z}\) 为元素族 \((\mathbf{X}_i\mathbf{Z})_{i \in I}\) （属于多项式环 \(\mathbf{A}[\mathbf{X}][\mathbf{Z}]\) ）。 \(\mathbf{Z}^k\) 在 \(u(\mathbf{X}, \mathbf{Z})\) 中的系数是次数为 \(k\) 的、 \(u\) 的齐次分量（对每个正整数 \(k\) ）。

只需在 \(u\) 是单项式的情形下证明此命题，而在此情形下结论是显然的。

推论. --- 对于一个多项式 \(u \in A[X]\) 为次数 \(k\) 的齐次多项式，其充分必要条件为：

\[
u (\mathbf {X} \cdot \mathbf {Z}) = u (\mathbf {X}) \cdot \mathbf {Z} ^ {k}.
\]

注记。——设 \(x \in A^{I}\) ，并设 f 为映射 \(u \mapsto u(x)\) （从 A{[}X{]} 到 A 中）。给定一个 A-模 M ，我们考虑同态 \(1 \otimes f\) （从 M{[}X{]} = M \(\otimes_{A} A[X]\) 到 \(\mathbf{M} \otimes_{\mathbf{A}} \mathbf{A} = \mathbf{M}\) 中）。对每个 \(v \in \mathbf{M}[\mathbf{X}]\) ，我们有 \((1 \otimes f)(v) = v(\mathbf{x})\) 。若 \(v = \sum_{\nu \in \mathbf{N}^{(I)}} e_{\nu} X^{\nu}\) ，则 \(v(\mathbf{x}) = \sum_{\nu \in \mathbf{N}^{(I)}} \mathbf{x}^{\nu} e_{\nu}\) 。

